% Generated draft: requires mathematical and source review.
\documentclass[11pt]{article}
\usepackage[margin=25mm]{geometry}
\usepackage{fontspec}
\setmainfont{Latin Modern Roman}
\usepackage{amsmath,amssymb,mathtools,unicode-math}
\setmathfont{Latin Modern Math}
\usepackage{enumitem,needspace,xurl,hyperref}
\setlength{\parindent}{0pt}
\setlength{\parskip}{5pt}
\setlength{\emergencystretch}{3em}
\newcommand{\N}{\mathbb{N}}
\newcommand{\Z}{\mathbb{Z}}
\newcommand{\Q}{\mathbb{Q}}
\newcommand{\R}{\mathbb{R}}
\newcommand{\C}{\mathbb{C}}
\newcommand{\F}{\mathbb{F}}
\newcommand{\A}{\mathbb{A}}
\newcommand{\PP}{\mathbb{P}}
\newcommand{\E}{\mathbb{E}}
\newcommand{\eps}{\varepsilon}
\newcommand{\dd}{\,\mathrm{d}}
\begin{document}
\section*{UnsolvedMath ID: 30002830}
\textbf{Source problem code:} catalogue entry
\subsection{Definitions and mathematical statement}
Let \(k\) denote the base field intended in the source, let \(t_1,t_2\) be algebraically independent over \(k\), and put \(K=k(t_1,t_2)\). The supplied record does not specify \(k\), its characteristic, or the construction of the Ueno-type variety \(X_{4,6}\). An integral \(k\)-variety \(V\) of dimension \(d\) is \emph{\(k\)-rational} if it is birational over \(k\) to \(\PP^d_k\), equivalently if
\[
 k(V)\cong_k k(u_1,\ldots,u_d),
\]
where \(u_1,\ldots,u_d\) are algebraically independent over \(k\). Birational modifications preserve the function field over the specified base field.

The supplied statement describes \(X_{4,6}\) birationally by a diagonal cubic surface \(S\) over \(K\), but supplies neither its equation nor the precise birational presentation. On the interpretation that this description gives a function-field identification
\[
 k(X_{4,6})\cong_k K(S),
\]
the intended rationality target is
\[
 k(X_{4,6})\cong_k k(u_1,u_2,u_3,u_4),
\]
since a surface over \(K\) has function field of transcendence degree \(4\) over \(k\).

\textbf{Problem as supplied.} Can the cubic surface in this birational description be modified birationally, analogously to the conic-bundle treatment of \(X_{3,4}\), to prove that \(X_{4,6}\) is rational over the intended base field? The construction of \(X_{3,4}\), the referenced conic-bundle treatment, and the permitted modifications are not specified. Consequently, this is a conditional formulation of the target, pending recovery of the missing definitions and hypotheses from the cited report. Rationality of \(S\) over \(K\) would be sufficient under the displayed presentation, but the record does not require that stronger intermediate conclusion.
\subsection{Short English statement}
The question asks whether birational modifications of the diagonal cubic-surface description of \(X_{4,6}\), inspired by the conic-bundle treatment of \(X_{3,4}\), can prove rationality. The supplied record omits the base-field hypotheses, the definition of the varieties, and the cubic equation, so the formulation remains conditional on those source details.
\subsection{Sources}
\begin{itemize}
\item Local UnsolvedMath registry: supplied record with permanent ID \(30002830\), exact source problem code \texttt{catalogue entry}, in the Oberwolfach Reports Open Problems collection. The record associates the problem with the year 2015; no individual proposer is supplied.
\item \emph{Algebraic Geometry} (2015), Oberwolfach Reports, pp.\ 783--836, DOI \texttt{10.4171/owr/2015/15}. Supplied report reference: \url{https://doi.org/10.4171/owr/2015/15}. The registry labels its statement \texttt{corrected\_verified} and reports a check against the official report; that verification is supplied provenance and has not been independently repeated here.
\item Catanese--Oguiso--Verra, described in the supplied record as proving unirationality of \(X_{4,6}\): \url{https://arxiv.org/abs/1506.01925}.
\item Mellit, described in the supplied record as giving a rationality argument conditional on heuristic curve-counting conjectures: \url{https://doi.org/10.1080/10586458.2019.1691088}.
\item Literature assessment dated \texttt{2026-08-21}: the supplied curation labels the problem open and assesses rationality as unresolved, distinguishing the stated unirationality result from Mellit's conditional rationality argument. This is supplied dated curation, not an independently verified literature assessment; the record explicitly requests independent verification of time-sensitive status and cited claims.
\end{itemize}
\end{document}
